\section{A finite first-layer certificate and a uniform fixed cubic bound}
\label{sec:cubic-layer-certificate}

Let $E_1=\mathbb Q(\mu_9,\sqrt[3]{269})$.  The exhaustive
class-group certification attempted for this degree-eighteen field
was stopped: its unconditional prime-search bound was
$9173231824586$.  Its preliminary class-group and logarithmic-class
outputs are not used as certified answers.  Instead, the following
finite certificate uses proposed algebraic objects from that
calculation and checks their required properties directly.
The verification scripts and complete outputs are retained in
\texttt{evidence/cubic\_logarithmic/}.

\begin{theorem}[AC-K15: a certified first-layer upper bound]
\label{AC-K15}
The actual degree-zero logarithmic class group satisfies
\[
       \#\widetilde{\operatorname{Cl}}_3(E_1)\leq3^6.    \tag{KC.1}
\]
Together with AC-K13 and AC-K16--AC-K17, this proves that the
split unramified Iwasawa module $X'$ of the fixed field $E$ is
finite.  Consequently there is a fixed finite integer
\[
 C_E=1+\operatorname{length}_{\mathcal O}X'
\]
such that, for every $e\geq1$ and every integer $t\geq2$,
\[
 3\leq m_e\leq C_E,\qquad
 \sum_{1\leq e<t}m_e\leq C_E(t-1).                   \tag{KC.2}
\]
The sharper equality $m_e=4$ for $3\mid e$ remains valid.
All assertions are unconditional.  The proof certifies the
upper bound (KC.1), rather than asserting exact first-layer
class invariants from an uncertified class-group routine.
No numerical value for $C_E$ is claimed.
\end{theorem}
\begin{proof}
\proofstep{1. Four explicit unramified cubic characters.}
The field used in the computation is the exact compositum of
$Y^3-269$ and $Y^6+Y^3+1$.  The compositum output retains its
embeddings; its chosen primitive element is
$\zeta_9-\sqrt[3]{269}$.  The degree-eighteen polynomial and
embeddings are in \texttt{certify\_layer1.out}, and the polynomial
is repeated in \path{layer1_unramified_certificate.out}.
The test \texttt{nfcertify(b)} returned the empty vector,
certifying the maximal order used in subsequent ideal operations.
This is a maximal-order check, not a class-group certification.

The file \texttt{layer1\_classfields.bin} retains a number-field
structure, a provisional ray-class structure, and four explicit
polynomials
\[
                         X^3-a_i\quad(1\leq i\leq4).
\]
Their algebraic coefficients $a_i\in E_1$ are printed completely
in \path{layer1_unramified_certificate.out}.
The routine proposing these polynomials is not used as a proof
of their class-field interpretation.  Each polynomial was
independently factored over the actual number field: the exact
factorization returned one factor of degree three.
For each, \texttt{rnfdisc(b,polynomial)[1]} returned the unit
ideal.  Thus each defines a degree-three extension unramified
at every finite prime.  The base field is totally imaginary,
so there are no additional real-place conditions.
The polynomials are pure cubics and $\mu_3\subset E_1$;
therefore each extension is cyclic.  This last argument is
specific to $X^3-a_i$, not to an arbitrary irreducible cubic.
They consequently define four genuine unramified Artin
characters on the ideal class group, with values in
$\mu_3\simeq\mathbb F_3$.

\proofstep{2. A certified upper bound for the ordinary three-class group.}
Let $g_1,g_2,g_3,g_4$ be the four explicit ideal matrices retained
in the field data, and let $\mathfrak q$ be the first prime
above seven in the retained prime decomposition.  Its norm is
$343$.  Cubic residue symbols of the $a_i$, evaluated at the
prime factors of $g_1,g_2,g_3,\mathfrak q$, give the matrix
\[
 \begin{pmatrix}
 0&2&0&0\\
 0&0&2&0\\
 2&0&0&0\\
 0&0&0&2
 \end{pmatrix},\qquad \det=1\quad\hbox{in }\mathbb F_3.
                                                               \tag{KC.3}
\]
All these test primes avoid the radicals' zero divisors.
Explicitly, at a prime $\mathfrak l$ of norm $N\mathfrak l$,
the script compares
$a_i^{(N\mathfrak l-1)/3}\bmod\mathfrak l$
with $1,\zeta_3,\zeta_3^2$; for an ideal product it sums
these exponents with the ideal-factorization multiplicities.
This is the Frobenius action on a cube root of $a_i$.
The root $\zeta_3$ is the sixth power of the displayed
order-eighteen root of unity.  Thus the computation is a
finite-field calculation of the actual Artin characters.

The same script supplies algebraic elements, in compact
factored form, and verifies by exact ideal arithmetic
\[
 \mathfrak q=(b)g_1^{21}g_4^2,
 \qquad
 g_1^{63}=(b_1),\quad g_2^{21}=(b_2),\quad
             g_3^3=(b_3),\quad g_4^3=(b_4).            \tag{KC.4}
\]
Every equality in (KC.4) returned verification value one.
Their complete elements and ideal matrices are in
\path{layer1_generation_certificate.out}; the reproducible
script has the same basename and suffix \texttt{.gp}.
These are direct ideal equalities, not assertions that the
provisional class-group generators form a basis.

Put $A_1=\operatorname{Cl}(E_1)_{(3)}$.
AC-K16 gives $\dim_{\mathbb F_3}A_1/\lambda A_1=2$.
Since $3$ is a unit times $\lambda^2$, the filtration by
$\lambda$ shows $\dim_{\mathbb F_3}A_1/3A_1\leq4$.
The rank-four pairing (KC.3) forces the classes of the four
tested ideals to span this quotient.  By the first relation
in (KC.4), they belong to the subgroup generated by the
$g_i$.  Nakayama's lemma for the finite three-group therefore
shows that the three-primary projections of the $g_i$
generate all of $A_1$.  The remaining relations in (KC.4)
now give the unconditional bound
\[
                       \#A_1\leq3^{2+1+1+1}=3^5.     \tag{KC.5}
\]
No claim that its order is exactly $3^5$ is needed.

\proofstep{3. Removing the classes at three saves a factor of three.}
Let $S$ be the three primes of $E_1$ over three.
For the radicals $a_1,a_2,a_4$, the exact test
\texttt{nfislocalpower(b,pr,a,3)} returned zero at every
$\mathfrak p\in S$; the third radical returned one.
The complete $4\times3$ result is
\[
 \begin{pmatrix}0&0&0\\0&0&0\\1&1&1\\0&0&0\end{pmatrix}.
                                                               \tag{KC.6}
\]
This tests local cubeness, independently of any class-group
or unit-group completeness assumption.  The implementation
in PARI/GP~2.19.1, \texttt{src/basemath/bnflog.c}, function
\texttt{nfislocalpower\_i}, strips the local valuation and
uses the finite principal-unit group.  Here the local
ramification index is six and its test level is
$6\cdot3/(3-1)+1=10$.  The public command semantics are
recorded in the same release's manual under
\texttt{nfislocalpower}.  Thus zero in (KC.6) is a certified
local noncube test, not a floating-point approximation.

Choose the first extension and any prime in $S$.
It is unramified by Step~1 and nonsplit by (KC.6), so its
Frobenius has order three.  The corresponding prime class
has a nontrivial image in an unramified character, and its
three-primary projection has order at least three.
It follows from (KC.5) that the actual localized class group
satisfies
\[
 \#\operatorname{Cl}(\mathcal O_{E_1}[1/3])_{(3)}
                         \leq3^4.                    \tag{KC.7}
\]

\proofstep{4. Two actual units bound the logarithmic kernel.}
Each completion of $E_1$ at three is $\mathbb Q_3(\mu_9)$.
Its logarithmic ramification and residue indices are
$(\widetilde e,\widetilde f)=(2,3)$, also independently
returned for all three primes by \texttt{bnflogef}.
Choose logarithmic degree $d=3\log_3(4)$ at each of these
three places.  This is the normalization allowed by
Belabas--Jaulent, Corollary~3.5 and Remark~3.6, author's
pp.~8--9.  For a genuine $S$-unit $u$, put
\[
  \ell_i(u)=-\frac{\log_3 N_{(E_1)_{\mathfrak p_i}/\mathbb Q_3}(u)}
                         {3\log_3(4)}\in\mathbb Z_3.
\]
The product formula gives $\sum_i\ell_i(u)=0$ exactly.

The script \path{layer1_logkernel_certificate.gp} retains
eleven explicit proposed $S$-units.  It verifies, by exact
ideal factorization, that every prime in every one's divisor
lies above three.  This only proves they are actual $S$-units;
it does not assume they are a full basis.  For units numbered
five and ten, the first two rows of their logarithmic
valuation matrix are
\[
       \begin{pmatrix}39&51\\21&30\end{pmatrix}
                  \pmod{81}.                         \tag{KC.8}
\]
Its determinant is $18\pmod{81}$, and hence has three-adic
valuation exactly two.  The script computes the local norms
as resultants with the three degree-six factors of the field
polynomial over $\mathbb Q_3$.  It starts with factor precision
sixty; the reported matrix precision is at least fifty,
and the determinant precision is fifty-one.  These are
absolute three-adic precision bounds.  In particular the
four digits required for the displayed reduction modulo
$81$ are available with a large margin.  Independent
replays at precisions forty and seventy give determinant
valuation two again.  Permuting the local factors can
change the displayed minor by a sign; it does not change
its valuation or the lattice index.

The lattice of degree-zero divisors supported in $S$ is
\[
 D_S^0=\{(z_1,z_2,z_3)\in\mathbb Z_3^3:
                                z_1+z_2+z_3=0\}.
\]
Its first two coordinates identify it with $\mathbb Z_3^2$.
By (KC.8), the logarithmic divisors of just these two actual
$S$-units span a sublattice of index nine in $D_S^0$.
Quotienting by all principal logarithmic divisors can only
reduce this order.  Therefore the degree-zero subgroup of
logarithmic classes generated at three has order at most
nine.

\proofstep{5. The logarithmic-class upper bound.}
The logarithmic divisor-class presentation has the exact
sequence
\[
 0\longrightarrow\widetilde{\operatorname{Cl}}_3(E_1)(3)
 \longrightarrow\widetilde{\operatorname{Cl}}_3^{\rm all}(E_1)
 \longrightarrow\operatorname{Cl}(\mathcal O_{E_1}[1/3])_{(3)}
 \longrightarrow0,                                    \tag{KC.9}
\]
where the middle group here includes the degree direction
and the left group is generated by the primes over three.
The source is Belabas--Jaulent, beginning of \S3, author's
p.~6, with the $S$-unit relation presentation in \S3.3,
p.~10.  Upon restriction to degree zero, surjectivity at
right need not persist; the image is a subgroup of the
last group.  Its order is at most $3^4$ by (KC.7), while
the kernel has order at most $3^2$ by Step~4.
Thus the degree-zero middle group is finite and has order
at most $3^6$, proving (KC.1).  This argument explains why
neither a saturated unit basis nor a full certified class
presentation is needed: all extra true relations only
make the group smaller.

\proofstep{6. All positive free ranks are excluded.}
Use $T=X'_{\rm tors}$, $X'/T\simeq\mathcal O^r$ and $G$
from AC-K16, and put $P(z)=\det(zI-G)$.
Suppose $r>0$.  AC-K17 proves
\[
 P\in\mathbb Z_3[z],\quad
 P(z)\equiv(z-1)^r\pmod3,\quad v_\lambda P(1)=2,
 \quad T\simeq\mathbb F_3[B]/B^k,
                                                               \tag{KC.10}
\]
where $B=\gamma-1$, $k\geq1$, and $r+k\in\{6,7\}$.
For clarity, the torsion assertion follows by applying
Nakayama to $T/(3,B)T\simeq\mathbb F_3$, choosing a
complex-conjugation eigenvector as its generator, and using
$c\sigma c=\sigma^{-1}$.  Conjugation is scalar on the
resulting cyclic $\mathbb Z_3[[B]]$-module, which forces
$\sigma=1$ and hence $\lambda T=0$.

The weight-zero descent over $E_1$ is
\[
                X'/(\gamma^3-1)X'
                       =\widetilde{\operatorname{Cl}}_3(E_1).
\]
It is finite by (KC.1).  Thus $G^3-I$ has nonzero determinant,
and the snake-lemma length formula used in (KI.7) gives
\[
 \operatorname{length}_{\mathcal O}
                  \widetilde{\operatorname{Cl}}_3(E_1)
   =\min(k,3)+v_\lambda\det(G^3-I).                  \tag{KC.11}
\]
Indeed $\gamma^3-1=B^3$ on $T$, since $T$ has characteristic
three.  If $r=1$, (KC.10) gives $k\geq5$, while
$G=(g)$ with $g\in1+3\mathbb Z_3$ and $v_3(g-1)=1$.
Lifting the exponent gives $v_\lambda(g^3-1)=4$.
The right side of (KC.11) is therefore $3+4=7$.

If $r\geq2$, write $P(z)=(z-1)^r+3Q(z)$ with
$Q\in\mathbb Z_3[z]$.  Both $P(\omega)$ and
$P(\omega^2)$ belong to $3\mathcal O$, because
$\omega-1$ and $\omega^2-1$ have $\lambda$-valuation one
and $r\geq2$.  The factorization
\[
 \det(G^3-I)=\pm P(1)P(\omega)P(\omega^2)
\]
therefore has $\lambda$-valuation at least six.
The torsion contribution $\min(k,3)$ is at least one,
so (KC.11) again has length at least seven.
Both cases contradict (KC.1).  Consequently $r=0$.
Finite generation in AC-K16 now makes $X'$ finite.

Finally $W_e(E)$ is a quotient of this fixed finite module
by the twisted endomorphism $4^e\gamma-1$.
Its $\mathcal O$-length is at most that of $X'$.
Equations (KL.1)--(KL.2) prove (KC.2), and AC-K14 supplies
the equality on multiples of three.
\end{proof}

\paragraph{Scope of the fixed-field result.}
The coefficient-prime-three contribution for this one cubic
character now has a uniform linear budget in the number of
weights.  This is stronger than mere weight-by-weight
finiteness.  It does not yet estimate other coefficient
primes, the varying CM hosts in \texttt{SupplyChi}, or the
full sum required in H04.4.  The constant $C_E$ and the two
remaining classes of individual weights modulo three have
not been explicitly evaluated.
